\begin{proof}
Let $M_1,M_2,\ldots$ be an enumeration of probabilistic machines running in $2^{0.1n}$ time. Let $n_1$ be a large enough constant and define the sequence $\{n_i\}$ where $n_i=2^{n_{i-1}}$ for each $i\geq2$. We let $Q=\mathcal O$ itself be the dense property without pseudodeterministic construction algorithms. In particular, we want that for every $i\in\mathbb N$, $M_i(1^{n_i})$ fails to generate a canonical string $x\in\mathcal O\cap\{0,1\}^{n_i}$. Naturally, our oracle construction proceeds in stages, where for each $i\in\mathbb N$, the $n_i$-th slice of $\mathcal O$ is constructed by carefully diagonalizing against $M_i$; on the other hand, if $n\in\mathbb N$ is not in the sequence $\{n_i\}$, then we simply let $\mathcal O$ accept every string of length $n$.

Now we show how to construct the $n_i$-th slice of $\mathcal O$ in the $i$-th stage. The idea is simple: construct the (purportedly pseudodeterministic) decision tree $\mathcal T$ corresponding to $M_i$, find an exponential-length input $tt$ of $\mathsf{Find1}$ on which $\mathcal T$ fails, and let the truth table of the $n_i$-th slice of $\mathcal O$ be $tt$. More precisely:
\begin{itemize}
\item \textbf{Converting $M_i$ into a distribution of decision trees.} We define the distribution of decision trees $\mathcal T$. The input to $\mathsf{Find1}$ has length $2^{n_i}$ and is considered as the truth table of the $n_i$-th slice of $\mathcal O$. Note that since $2^{0.1n_i}\ll n_{i+1}$, for every $n'\leq2^{0.1n_i}$ such that $n'\neq n_i$, the $n'$-th slice of $\mathcal O$ are fixed. To sample a (deterministic) decision tree $T$ from $\mathcal T$, sample a sequence of random coins fed to $M_i$ and simulate $M_i$ on the oracle $\mathcal O$. Whenever $M_i$ makes a query $\mathcal O(x)$, if $|x|=n_i$, then the decision tree asks the $x$-th bit of our input; otherwise we return the already-fixed value of $\mathcal O(x)$.
\item \textbf{Invoking the lower bound.} Note that $\mathcal T$ only makes $2^{0.1n_i}<o(\sqrt{2^{n_i}})$ queries. By Theorem 51, there exists an input $tt\in\{0,1\}^{2^{n_i}}$ such that $\mathcal T$ fails to solve $\mathsf{Find1}$ on input $tt$ pseudodeterministically. That is, for every valid output $x$ such that $tt_x=1$, $\Pr[\mathcal T(tt)=x]<2/3$. Let the truth table of the $n_i$-th slice of $\mathcal O$ to be such a $tt$.
\end{itemize}
It is easy to see that for every probabilistic machine $A$ that runs in $2^{o(n)}$ time with oracle access to $\mathcal O$, and every infinite sequence of outputs $\{x_n\}_{n\in\mathbb N}$ where each $x_n\in\mathcal O\cap\{0,1\}^n$, for every $i$ such that $A=M_i$, the probability that $A(1^{n_i})$ outputs $x_{n_i}$ is at most $3/4$.

\textbf{Complexity of $\mathcal O$.} Let $\mathcal O_{\leq n_{i-1}}$ denote the description of the oracle $\mathcal O$ up to the $n_{i-1}$-th slice (note that $\mathcal O_{\leq n_{i-1}}$ can be described in $\operatorname{poly}(2^{n_{i-1}})\leq\operatorname{poly}(n_i)$ bits), and $x\in\{0,1\}^{n_i}$. Note that $\mathcal O_{\leq n_{i-1}}$ and $tt$ together defines the oracle $\mathcal O$ up to input length $n_{i+1}-1$, so the behavior of $M_i^{\mathcal O}$ is completely determined by $M_i$, $\mathcal O_{\leq n_{i-1}}$, and $tt$. Let $\mathsf{ProbEst}(\mathcal O_{\leq n_{i-1}},tt,M_i,x)$ be an oracle such that
\begin{align*}
\Pr[M_i^{\mathcal O}(1^{n_i})=x]>3/4&\Longrightarrow\mathsf{ProbEst}(\mathcal O_{\leq n_{i-1}},tt,M_i,x)=1,\\
\Pr[M_i^{\mathcal O}(1^{n_i})=x]<2/3&\Longrightarrow\mathsf{ProbEst}(\mathcal O_{\leq n_{i-1}},tt,M_i,x)=0.
\end{align*}
By Theorem 23, the oracle $\mathsf{ProbEst}$ can be implemented in $\mathsf{DTIME}[2^{O(n_i)}]^{\mathsf{PH}}$. Consider the following algorithm with oracle access to $\mathsf{ProbEst}$ that prints a truth table $tt$. The algorithm maintains the prefix of a truth table which is initially the empty string, and extends this prefix bit by bit. Suppose that we have a prefix of length $\ell$, denoted as $tt'\in\{0,1\}^{\ell}$. To fix the $(\ell+1)$-st bit of $tt'$, we check if there exists a truth table $tt$ such that (1) for every $x\in\{0,1\}^{n_i}$, $\mathsf{ProbEst}(\mathcal O_{\leq n_{i-1}},tt,M_i,x)=0$, and (2) $tt'\circ0$ ($tt'$ concatenated with a bit $0$) is a prefix of $tt$. If there is such a $tt$, then we append $0$ to the end of $tt'$; otherwise we append $1$ to the end of $tt'$.

It is clear that the algorithm always outputs a truth table $tt$ on which $M_i$ fails. Since $\mathsf{ProbEst}$ runs in $\mathsf{DTIME}[2^{O(n_i)}]^{\mathsf{PH}}$, the whole algorithm also runs in $\mathsf{DTIME}[2^{O(n_i)}]^{\mathsf{PH}}$. Using this algorithm, it is easy to construct the oracle $\mathcal O$ in $\mathsf{EXPH}$.
\end{proof}
